\documentclass[aspectratio=169,11pt]{beamer}
\usepackage{paperdeck}
\title[Gemma 2]{Gemma 2: Improving Open Language Models at a Practical Size}
\subtitle{Gemma Team, Google DeepMind (2024), arXiv:2408.00118}
\author[Paper summary]{Paper summary for the CS-781 capstone}
\date{}
\begin{document}

\begin{frame}\titlepage\end{frame}

\paperglance
{Gemma Team, Google DeepMind (M.\ Riviere, S.\ Pathak, P.\ G.\ Sessa, C.\ Hardin, et al.). \emph{Gemma 2: Improving Open Language Models at a Practical Size}. Technical report, arXiv:2408.00118, 2024.}
{Open decoder-only models (2B, 9B, 27B) that combine interleaved local/global attention, grouped-query attention and other refinements. The 2B and 9B models are trained by \textbf{knowledge distillation} instead of plain next-token prediction.}
{Our A4 follows the same recipe at small scale: local + global layers, GQA, RoPE, RMSNorm and a gated FFN. Distillation is mandatory in our Phases 2--7.}
{Main deck: ``A4: efficient decoder'' and ``Compression: knowledge distillation''.}

\begin{frame}{Roadmap}
  \begin{multicols}{2}\small\tableofcontents\end{multicols}
\end{frame}

% ------------------------------------------------------------------
\section{The problem}

\begin{frame}{Small models are under-trained}
  \begin{columns}[T]
    \begin{column}{0.52\textwidth}
      \small
      \begin{itemize}
        \item Recent small models improved mainly by \textbf{training on more tokens}.
        \item But gains grow only logarithmically with data: the latest small models needed up to 15T tokens to improve by less than 1--2\%.
        \item So small models are still under-trained, and more raw text is an expensive fix.
      \end{itemize}
      \vspace{1mm}
      \textbf{Question}: can we give a small model \emph{more information per training step} instead of more steps?
    \end{column}
    \begin{column}{0.44\textwidth}
      \begin{exampleblock}{The answer: distillation}
        \small
        Replace the one-hot next-token target with the \textbf{full probability distribution} from a larger teacher model.
        \vspace{1mm}

        The report uses this to \emph{simulate training beyond the available data}: the 2B and 9B students see more than 50$\times$ the compute-optimal number of tokens.
      \end{exampleblock}
    \end{column}
  \end{columns}
\end{frame}

% ------------------------------------------------------------------
\section{The basic idea: architecture}

\begin{frame}{Model overview}
  \centering
  \footnotesize
  \begin{tabular}{@{}lccc@{}}
        \toprule
        & 2B & 9B & 27B\\
        \midrule
        $d_{\text{model}}$ & 2304 & 3584 & 4608\\
        Layers & 26 & 42 & 46\\
        Query heads / KV heads & 8 / 4 & 16 / 8 & 32 / 16\\
        Head size & 256 & 256 & 128\\
        FFN (GeGLU) dim & 18,432 & 28,672 & 73,728\\
        Sliding window / global span & \multicolumn{3}{c}{4,096 / 8,192}\\
        Vocabulary & \multicolumn{3}{c}{256,128 (tied embeddings)}\\
        Training tokens & 2T & 8T & 13T\\
        Trained by & distillation & distillation & from scratch\\
        \bottomrule
      \end{tabular}
  \vspace{2mm}

  \raggedright
  \begin{columns}[T]
    \begin{column}{0.48\textwidth}
      \begin{block}{Shared with Gemma 1}
        \small
        Decoder-only Transformer, 8,192-token context, RoPE positions, approximated GeGLU, tied embeddings.
      \end{block}
    \end{column}
    \begin{column}{0.48\textwidth}
      \begin{exampleblock}{Embeddings are costly here}
        \small
        The 256k vocabulary gives the 2B model 590M embedding parameters (2.02B non-embedding). Our 32k vocabulary costs 24.6M of 82M.
      \end{exampleblock}
    \end{column}
  \end{columns}
\end{frame}

\begin{frame}{Interleaved local and global attention}
  \begin{columns}[T]
    \begin{column}{0.52\textwidth}
      \small
      Layers alternate: \textbf{local} sliding-window attention (4,096 tokens), then \textbf{global} attention (8,192 tokens), every other layer.
      \begin{itemize}
        \item Local layers are cheap: cost and KV-cache grow with the window, not the full length.
        \item Global layers keep long-range information flowing.
      \end{itemize}
      \vspace{1mm}
      \textbf{Ablation}: changing the 9B model's window \emph{only at inference}:
      \vspace{1mm}

      \begin{tabular}{@{}lccc@{}}
        \toprule
        Sliding window & 4096 & 2048 & 1024\\
        \midrule
        Validation perplexity & 1.63 & 1.63 & 1.64\\
        \bottomrule
      \end{tabular}
    \end{column}
    \begin{column}{0.44\textwidth}
      \begin{exampleblock}{Reading it}
        \small
        Shrinking the local window 4$\times$ at inference barely changes perplexity. The window can be a knob for a small speed gain.
      \end{exampleblock}
      \begin{block}{Our A4 at small scale}
        \small
        A4 uses a 256-token window in local layers and full attention in every \emph{third} layer (Gemma 2: every second). Its context is 1,024 tokens.
      \end{block}
    \end{column}
  \end{columns}
\end{frame}

\begin{frame}{Logit soft-capping}
  \begin{columns}[T]
    \begin{column}{0.52\textwidth}
      \small
      Keep logits within $(-c,\,c)$ with a smooth squash:
      \[
        \text{logits}\leftarrow c\cdot\tanh(\text{logits}/c)
      \]
      \begin{itemize}
        \item $c=50$ for attention scores in every layer.
        \item $c=30$ for the final output logits.
      \end{itemize}
      Small values pass almost unchanged; very large values are squeezed toward $\pm c$. This prevents a few extreme logits from destabilising training.
    \end{column}
    \begin{column}{0.44\textwidth}
      \begin{exampleblock}{Easy example ($c=30$)}
        \small
        \begin{tabular}{@{}lc@{}}
          \toprule
          Raw logit & Capped\\
          \midrule
          5 & $30\tanh(0.167)=4.95$\\
          10 & $30\tanh(0.333)=9.65$\\
          30 & $30\tanh(1)=22.85$\\
          60 & $30\tanh(2)=28.9$\\
          300 & $\approx30.0$\\
          \bottomrule
        \end{tabular}
        \vspace{1mm}

        Ordinary logits barely move; a runaway logit of 300 is held at 30.
      \end{exampleblock}
    \end{column}
  \end{columns}
\end{frame}

\begin{frame}{Normalisation and grouped-query attention}
  \begin{columns}[T]
    \begin{column}{0.5\textwidth}
      \begin{block}{Pre-norm \emph{and} post-norm with RMSNorm}
        \small
        RMSNorm is applied to both the \textbf{input and the output} of each sub-layer (attention and FFN), for training stability:
        \[
          x\leftarrow x+\text{RMSNorm}\big(f(\text{RMSNorm}(x))\big)
        \]
        {\scriptsize Our A4 uses only pre-norm.}
      \end{block}
    \end{column}
    \begin{column}{0.46\textwidth}
      \begin{block}{Grouped-query attention}
        \small
        Two query heads share each KV head (\code{num\_groups = 2}).
        \vspace{1mm}

        \textbf{Ablation (9B)}, average over 4 benchmarks:
        \vspace{1mm}

        \begin{tabular}{@{}lc@{}}
          \toprule
          MHA & 50.3\\
          GQA & 50.8\\
          \bottomrule
        \end{tabular}
        \vspace{1mm}

        Little change in quality, with fewer parameters and faster inference.
      \end{block}
    \end{column}
  \end{columns}
\end{frame}

% ------------------------------------------------------------------
\section{The basic idea: training by distillation}

\begin{frame}{Knowledge distillation as the pre-training objective}
  \begin{columns}[T]
    \begin{column}{0.52\textwidth}
      \small
      For each position, the teacher gives a distribution $P_T(x\mid x_c)$ over the next token. The student minimises the cross-entropy to it:
      \[
        \min_{P_S}\;\sum_x -P_T(x\mid x_c)\,\log P_S(x\mid x_c)
      \]
      \begin{itemize}
        \item Standard training uses a one-hot target: only the true next token counts.
        \item Distillation gives a \emph{richer gradient}: every plausible token gets credit.
        \item Used for the 2B and 9B models. The 27B model was trained from scratch.
      \end{itemize}
    \end{column}
    \begin{column}{0.44\textwidth}
      \begin{exampleblock}{Easy legal example}
        \small
        Context: ``The appeal is hereby \_\_''. True next word: ``dismissed''.
        \vspace{1mm}

        One-hot target: dismissed 1.0.\\
        Teacher target: dismissed 0.62, rejected 0.20, allowed 0.10, disposed 0.05, \ldots
        \vspace{1mm}

        From one example the student learns four useful facts instead of one, much as in Hinton et al.\ (2015).
      \end{exampleblock}
    \end{column}
  \end{columns}
\end{frame}

\begin{frame}{Does distillation really help? Ablations}
  \begin{columns}[T]
    \begin{column}{0.5\textwidth}
      \small
      \textbf{2B model, 500B tokens}, from scratch vs distilled from a 7B teacher:
      \vspace{1mm}

      \begin{tabular}{@{}lc@{}}
        \toprule
        & Average (3 benchmarks)\\
        \midrule
        From scratch & 60.3\\
        Distilled & \textbf{67.7}\\
        \bottomrule
      \end{tabular}
      \vspace{2mm}

      \textbf{Across student sizes} (validation perplexity, 7B teacher):
      \vspace{1mm}

      \begin{tabular}{@{}lccc@{}}
        \toprule
        & 200M & 400M & 1B\\
        \midrule
        From scratch & 23 & 19 & 17\\
        Distilled & \textbf{21} & \textbf{17} & \textbf{15}\\
        \bottomrule
      \end{tabular}
    \end{column}
    \begin{column}{0.46\textwidth}
      \begin{exampleblock}{Reading it}
        \small
        \begin{itemize}
          \item A 7.4-point average gain at the same data and size.
          \item The gain holds as the student grows from 200M to 1B.
          \item 200M is close to our size range (50--300M), so the result is directly relevant to us.
        \end{itemize}
      \end{exampleblock}
      \begin{block}{Wide vs deep (9B)}
        \small
        At equal parameters, a \textbf{deeper} network scored 52.0 vs 50.8 for a wider one: a small but consistent gain.
      \end{block}
    \end{column}
  \end{columns}
\end{frame}

% ------------------------------------------------------------------
\section{Results}

\begin{frame}{Pre-trained model results}
  \centering
  \footnotesize
  \begin{tabular}{@{}lccccc@{}}
        \toprule
        Benchmark & Gemma-1 2B & Gemma-2 2B & Gemma-1 7B & Gemma-2 9B & Gemma-2 27B\\
        \midrule
        MMLU & 42.3 & 52.2 & 64.4 & 71.3 & 75.2\\
        GSM8K & 15.1 & 24.3 & 51.8 & 68.6 & 74.0\\
        ARC-C & 48.5 & 55.7 & 61.1 & 68.4 & 71.4\\
        TriviaQA & 53.2 & 60.4 & 63.4 & 76.6 & 83.7\\
        HumanEval & 22.0 & 20.1 & 32.3 & 40.2 & 51.8\\
        \midrule
        Average (8) & 44.0 & 50.0 & 62.4 & 70.2 & 74.4\\
        \bottomrule
      \end{tabular}
  \vspace{1mm}

  {\scriptsize For comparison, Mistral 7B averages 61.0 and LLaMA-3 8B 61.9 on the same 8 benchmarks.}
  \vspace{1mm}

  \raggedright
  \begin{exampleblock}{Reading it}
    \small
    Gemma-2 9B beats LLaMA-3 8B by about 8 points on average. Gemma-2 2B beats Gemma-1 2B while trained on \emph{fewer} tokens (2T vs 3T); the report attributes the gain mainly to distillation, though the architecture changed too. The report claims the models compete with models 2--3$\times$ larger.
  \end{exampleblock}
\end{frame}

\begin{frame}{Post-training, robustness and safety}
  \begin{columns}[T]
    \begin{column}{0.5\textwidth}
      \small
      \textbf{Instruction tuning}:
      \begin{itemize}
        \item supervised fine-tuning on synthetic and human prompts, with responses mostly written by a larger \emph{teacher};
        \item on-policy distillation from the teacher;
        \item RLHF with a reward model much larger than the policy;
        \item \textbf{model merging}: averaging models trained with different hyperparameters.
      \end{itemize}
    \end{column}
    \begin{column}{0.46\textwidth}
      \begin{block}{Format robustness}
        \small
        Standard deviation of MMLU over 12 prompt and evaluation formats:
        \vspace{1mm}

        \begin{tabular}{@{}lc@{}}
          \toprule
          Gemma 2 2B / 9B / 27B & 2.1 / 0.9 / 1.0\\
          Mistral 7B & 6.9\\
          \bottomrule
        \end{tabular}
      \end{block}
      \begin{alertblock}{Responsible release}
        \small
        Data filtering for personal and unsafe content, evaluation-set decontamination, and safety evaluations. Users are told to run their own safety tests before deployment.
      \end{alertblock}
    \end{column}
  \end{columns}
\end{frame}

% ------------------------------------------------------------------
\section{Assessment and outlook}

\begin{frame}{Lessons for our capstone}
  \small
  \begin{tabularx}{\textwidth}{@{}lX@{}}
    \toprule
    Gemma 2 finding & What it suggests for us\\
    \midrule
    Distillation beats from-scratch for small students & Distil A4-small from A1 in Phase 1; expect gains in the mandatory Phase 2--7 distillation too\\
    Gain holds from 200M to 1B & Our 47--300M range is in scope\\
    Local window can shrink at inference & Try A4 with windows of 128/256/512 at generation time\\
    GQA $\approx$ MHA in quality & A4's 4 KV heads should cost little quality\\
    Deeper beats wider & At fixed parameters, test more layers with a smaller $d_{\text{model}}$\\
    Soft-capping stabilises training & A cheap addition if fp16 training on a T4 diverges\\
    Model merging helps & Average checkpoints from different runs before submission\\
    \bottomrule
  \end{tabularx}
\end{frame}

\begin{frame}{Strengths and limitations}
  \begin{columns}[T]
    \begin{column}{0.48\textwidth}
      \begin{exampleblock}{Strengths}
        \small
        \begin{itemize}
          \item Strong open models at practical sizes.
          \item Clear ablations for distillation, GQA, depth and window size.
          \item Shows distillation as a way to go beyond available data.
          \item Combines several proven techniques in one design.
          \item Open weights for the community.
        \end{itemize}
      \end{exampleblock}
    \end{column}
    \begin{column}{0.48\textwidth}
      \begin{alertblock}{Limitations}
        \small
        \begin{itemize}
          \item A technical report: training data and many details are not public.
          \item Ablations use few benchmarks, often averages over 3--4.
          \item Needs a strong, expensive teacher, and teacher inference over trillions of tokens.
          \item Mostly English; not built for strong multilingual use.
          \item The authors list factuality, adversarial robustness, reasoning and alignment as open problems.
        \end{itemize}
      \end{alertblock}
    \end{column}
  \end{columns}
\end{frame}

\begin{frame}{Future research directions}
  \begin{columns}[T]
    \begin{column}{0.5\textwidth}
      \small
      \textbf{Stated by the authors}
      \begin{itemize}
        \item Improve \textbf{factuality}.
        \item Improve \textbf{robustness to adversarial attacks}.
        \item Improve \textbf{reasoning}.
        \item Improve \textbf{alignment}.
      \end{itemize}
      \vspace{1mm}
      \textbf{Natural follow-ups}
      \begin{itemize}
        \item How small can the teacher--student gap be before distillation stops paying off?
        \item Cheaper distillation: store only the teacher's top-$k$ probabilities.
      \end{itemize}
    \end{column}
    \begin{column}{0.46\textwidth}
      \begin{block}{Open for legal NLP}
        \small
        \begin{itemize}
          \item A law firm could distil a large in-house model into a laptop-sized one, keeping client data private: the motivation of this capstone.
          \item Local attention suits long judgments, but party names and statutes from page 1 must reach page 40. How many global layers are enough?
          \item Can distillation from a general teacher transfer legal knowledge, or must the teacher itself be legal?
        \end{itemize}
      \end{block}
    \end{column}
  \end{columns}
\end{frame}

\begin{frame}{Summary}
  \begin{itemize}
    \item Small models are under-trained, and more raw data gives only logarithmic gains.
    \item \textbf{Gemma 2}: decoder-only, with interleaved local (4,096) and global (8,192) attention, GQA, RoPE, GeGLU, pre+post RMSNorm and logit soft-capping.
    \item The 2B and 9B models are trained by \textbf{distillation} from a larger teacher, over 50$\times$ the compute-optimal token count.
    \item Ablations: distillation 67.7 vs 60.3 from scratch; gains hold from 200M to 1B; GQA $\approx$ MHA; deeper beats wider; the window can shrink at inference.
    \item Gemma-2 9B averages 70.2 vs 61.9 for LLaMA-3 8B.
    \item For us: a large-scale version of A4's design, and strong evidence for our distillation experiments.
  \end{itemize}
\end{frame}

\begin{frame}{Reference}
  \small
  Gemma Team, Google DeepMind (2024). Gemma 2: Improving Open Language Models at a Practical Size. arXiv:2408.00118.
  \vspace{3mm}

  {\footnotesize Other work cited: Hinton et al.\ (2015) Distillation; Beltagy et al.\ (2020) Longformer (local attention); Shazeer (2020) GLU Variants; Su et al.\ (2021) RoPE; Hoffmann et al.\ (2022) Compute-optimal training; Ainslie et al.\ (2023) GQA; Agarwal et al.\ (2024) and Gu et al.\ (2024) on-policy distillation; Ramé et al.\ (2024) model merging.}
\end{frame}

\end{document}
